\expandafter\def\csname HistoricalNoteText4\endcsname{In einem Brief an \textsc{Leibnitz}. Horrebowii Opera T. II p. 142.}
\expandafter\def\csname HistoricalNoteText5\endcsname{Dieser ist also gleich dem Sinus der mittleren Horizontalparallaxe der Sonne, wenn die mittlere
Entfernung der Erde von der Sonne als Einheit gewählt wird.}
\clearpage\SourceLeaf{99}%
\pdfbookmark[1]{Einige Bemerkungen zur Vereinfachung der Berechnung geozentrischer Planetenpositionen}{bemerkungen1804}
\begin{center}
{\small EINIGE BEMERKUNGEN ZUR VEREINFACHUNG DER BERECHNUNG
DER GEOZENTRISCHEN POSITIONEN DER PLANETEN.\par}
\medskip\hrule\medskip
{\small Monatliche Correspondenz zur Beförderung der Erd- und Himmels-Kunde,\par
herausgegeben vom Freiherrn \textsc{von Zach}. Mai 1804.\par}
\medskip\hrule
\end{center}\medskip
Seit der Erfindung der Pendeluhren beziehen sich alle unsere Beobachtungen
der Fixsterne, Planeten und Kometen unmittelbar auf ihre Lage zum Äquator
statt zur Ekliptik. Auch unseren neuesten und besten Sternverzeichnissen und
Sternkarten liegen Rektaszension und Deklination statt Länge und Breite zugrunde.
Daher gibt es sehr häufig Anlass, aus den heliozentrischen Positionen der Planeten
und Kometen in ihren Bahnen ihre geozentrischen Positionen in Bezug auf den
Äquator zu berechnen. Noch häufiger ergäbe sich dieser Anlass, wenn man sich
entschlösse, in den astronomischen Ephemeriden durchgängig statt der wenig
nützlichen Längen und Breiten der Planeten die in jeder praktischen Hinsicht
viel brauchbareren Rektaszensionen und Deklinationen anzusetzen. Das hat der
vortreffliche \textsc{Römer} schon vor hundert Jahren empfohlen\HistoricalNote{4},
und es wird besonders für die beiden zuletzt entdeckten Planeten ganz
unentbehrlich, die so schwer zu beobachten und nur mithilfe sehr detaillierter
Sternkarten unter den kleinen Fixsternen in ihrer Umgebung auszumachen sind.
Ebenso oft würde die allgemeinere Befolgung eines anderen Vorschlags Anlass zu
jener Rechnung geben. Beim Vergleich der beobachteten Position
\SourceLeaf{100}%
eines Planeten oder Kometen mit seiner berechneten Position soll man die
beobachtete Rektaszension und Deklination unmittelbar zugrunde legen, statt
daraus, wie es gewöhnlich geschieht, erst eine sogenannte beobachtete Länge
und Breite abzuleiten. Die Vorteile dieses Verfahrens hat ein sachkundiger
Beurteiler bereits in Band V der Mon. Corr., S. 594, erwähnt.

Unter diesem Gesichtspunkt sind die geozentrische Länge und Breite des
Planeten lediglich Zwischengrößen zur Bestimmung seiner Lage zum Äquator.
Für die obigen Vorschläge spricht daher vielleicht zusätzlich, dass man sich
diese Zwischenrechnung und sogar die Reduktion der heliozentrischen Position
in der Bahn auf die heliozentrische Position in Bezug auf die Ekliptik ganz
ersparen kann. Aus der Position in der Bahn lassen sich die geozentrische
Rektaszension und Deklination unmittelbar mit sehr einfachen und flexiblen
Formeln ableiten, die in diesem Aufsatz entwickelt werden sollen. Zu diesen
Vorteilen kommt die große Leichtigkeit, mit der sich bei diesem Verfahren die
\textit{Parallaxe} auch dann berücksichtigen lässt, wenn sich der Planet außerhalb
des Meridians des Beobachtungsorts befindet. Das ist zwar seltener nötig, bei
anderen Methoden dann aber wesentlich umständlicher.

Man lege durch den Mittelpunkt der Sonne drei zueinander senkrechte Ebenen,
die erste parallel zum Erdäquator, die zweite durch die Äquinoktialpunkte
und die dritte somit durch die Sonnenwendepunkte. Die senkrechten Abstände
des Erdmittelpunkts von diesen Ebenen seien der Reihe nach $Z$, $Y$, $X$,
die Abstände eines Planeten von denselben Ebenen $z$, $y$, $x$. Diese Abstände
seien auf der Seite der ersten Ebene positiv, auf der der Nordpol liegt, bei
der zweiten auf der Seite der Sommersonnenwende und bei der dritten auf der
Seite der Frühlingstagundnachtgleiche. Demnach sind $z-Z$, $y-Y$, $x-X$ die auf
dieselbe Weise bestimmten senkrechten Abstände des Planeten von drei zu den
obigen Ebenen parallelen Ebenen durch den Erdmittelpunkt. Bezeichnet man also
die geozentrische Rektaszension des Planeten mit $\alpha$, seine Deklination mit
$\delta$ und seinen Abstand von der Erde mit $\Delta$, so ergibt sich
\[
\begin{aligned}
x-X&=\Delta\cos\delta\cos\alpha\\
y-Y&=\Delta\cos\delta\sin\alpha\\
z-Z&=\Delta\sin\delta
\end{aligned}
\]
\SourceLeaf{101}%
Man findet also $\alpha$ mit der Formel $\operatorname{tang}\alpha=\frac{y-Y}{x-X}$.
Das positive oder negative Vorzeichen des Zählers muss darüber entscheiden,
ob $\alpha$ in den ersten beiden oder den letzten beiden Quadranten anzusetzen ist.
Dann erhält man
\[
\Delta\cos\delta=\frac{x-X}{\cos\alpha}=\frac{y-Y}{\sin\alpha},\quad
\text{und}\quad\operatorname{tang}\delta=\frac{z-Z}{\Delta\cos\delta}
\]

So erhält man die Rektaszension und Deklination des Planeten vom Erdmittelpunkt
aus gesehen. Sucht man sie so, wie sie von einem Punkt auf der Erdoberfläche aus
erscheinen, muss man in den obigen Formeln nur statt der Koordinaten $X$, $Y$, $Z$
des Erdmittelpunkts die Abstände des Beobachtungsorts von den drei Grundebenen
verwenden. Der Erdradius sei $=\rho$\HistoricalNote{5}, die Polhöhe des Beobachtungsorts
$=\varphi$ und die dortige Sternzeit im Augenblick der Beobachtung, im Bogenmaß
angegeben, oder die Rektaszension des kulminierenden Äquatorpunkts $=\vartheta$.
Dann ergeben sich, wie man leicht sieht, die Abstände
\[
\begin{gathered}
X+\rho\cos\varphi\cos\vartheta\\
Y+\rho\cos\varphi\sin\vartheta\\
Z+\rho\sin\varphi
\end{gathered}
\]
Dabei wird die Erde als Kugel angenommen. Sollte es nötig sein, auch ihre
sphäroidische Form zu berücksichtigen, wie es bei Kometen der Fall sein könnte,
die der Erde sehr nahe kommen, müsste man für $\rho$ nur den Abstand des
Beobachtungsorts vom Erdmittelpunkt und für $\varphi$ seine sogenannte
verbesserte Polhöhe einsetzen. Beide werden nach bekannten Regeln bestimmt.

Nun sieht man, dass nur noch eine bequeme Methode zur Bestimmung der Koordinaten
$X$, $Y$, $Z$, $x$, $y$, $z$ zu finden ist. Dazu beschreibe man eine Kugelfläche
mit unbestimmtem Radius um die Sonne. Auf ihr bezeichne $P$ den Nordpol der
Ekliptik, $p$ den Nordpol der Ebene der Planetenbahn, $K$ die Position der Erde,
$k$ die heliozentrische Position des Planeten und schließlich $\mathfrak{x}$,
$\mathfrak{y}$, $\mathfrak{z}$ die Pole der drei Grundebenen auf den Seiten, auf
denen die Abstände $x$, $y$, $z$ positiv angesetzt werden.
\SourceLeaf{102}%
Damit bezeichnet $\mathfrak{z}$ den Nordpol des Äquators, $\mathfrak{x}$ den
Frühlingspunkt und $\mathfrak{y}$ den Punkt des Äquators mit der Rektaszension
$90^0$. Wer eine Figur benötigt, kann sie danach leicht selbst zeichnen. Setzen
wir nun den Abstand der Erde von der Sonne $=R$, so ergibt sich offensichtlich
\[
\begin{aligned}
X&=R\cos\mathfrak{x}K\\
Y&=R\cos\mathfrak{y}K\\
Z&=R\cos\mathfrak{z}K
\end{aligned}
\]
Da im sphärischen Dreieck $\mathfrak{x}PK$ die Seite $PK=90^0$ ist und somit
$\cos\mathfrak{x}K=\sin\mathfrak{x}P\cos\mathfrak{x}PK$ gilt, folgt
\[
\begin{aligned}
X&=R\sin\mathfrak{x}P\cos\mathfrak{x}PK,\quad\text{und ebenso}\\
Y&=R\sin\mathfrak{y}P\cos\mathfrak{y}PK\quad\text{und}\\
Z&=R\sin\mathfrak{z}P\cos\mathfrak{z}PK
\end{aligned}
\]

Ganz entsprechend erhält man für die Koordinaten des Planeten, wenn wir
seinen Abstand von der Sonne mit $r$ bezeichnen,
\[
\begin{aligned}
x&=r\sin\mathfrak{x}p\cos\mathfrak{x}pk\\
y&=r\sin\mathfrak{y}p\cos\mathfrak{y}pk\\
z&=r\sin\mathfrak{z}p\cos\mathfrak{z}pk
\end{aligned}
\]

Wir legen hier ein für alle Mal fest, dass der sphärische Winkel
$\mathfrak{x}PK$ in dem Sinn zu verstehen ist, in dem der Schenkel $PK$ dem
Schenkel $P\mathfrak{x}$ \textit{in der Reihenfolge der Zeichen folgt}.
Er ist also von $KP\mathfrak{x}$ zu unterscheiden, und beide Winkel ergänzen
sich zu $360^0$. Ebenso ist jeder andere sphärische Winkel zu verstehen.
Mit dieser genaueren Festlegung können wir die Grundformeln der sphärischen
Trigonometrie ohne Weiteres auf Dreiecke mit Winkeln über $180^0$ ausdehnen.
Dadurch vermeiden wir die sonst nötige Unterscheidung mehrerer Einzelfälle.
Winkel, deren Differenz $360^0$ oder ein Vielfaches davon beträgt, gelten
durchweg als gleichbedeutend.

Wir betrachten nun zunächst die Koordinaten $X$, $Y$, $Z$ und setzen die
Schiefe der Ekliptik $=\varepsilon$ und die heliozentrische Länge der Erde
$=\lambda$ ($=$ geozentrische Länge der Sonne $+180^0$). In den obigen Formeln
gilt also $\mathfrak{x}P=90^0$, $\mathfrak{y}P=90^0+\varepsilon$,
$\mathfrak{z}P=\varepsilon$, $\mathfrak{x}PK=\lambda$,
$\mathfrak{y}PK=\mathfrak{z}PK=\lambda-90^0$, und folglich
\SourceLeaf{103}%
\[
\begin{aligned}
X&=R\cos\lambda\\
Y&=R\sin\lambda\cos\varepsilon\\
Z&=R\sin\lambda\sin\varepsilon
\end{aligned}
\]

Für den Planeten setzen wir der Kürze halber $\mathfrak{x}p=a$,
$\mathfrak{y}p=b$, $\mathfrak{z}p=c$, seinen Abstand längs der Bahn vom
aufsteigenden Knoten auf der Ekliptik $=t$ und die Winkel $\mathfrak{x}pP$,
$\mathfrak{y}pP$, $\mathfrak{z}pP$ der Reihe nach $=A$, $B$, $C$. Man sieht
leicht, dass $Ppk=t-90^0$ (oder nach der obigen Bemerkung $=t+270^0$) gilt,
also $\mathfrak{x}pk=A+t-90^0$, $\mathfrak{y}pk=B+t-90^0$,
$\mathfrak{z}pk=C+t-90^0$. Damit erhält man
\[
\begin{aligned}
x&=r\sin a\sin(A+t)\\
y&=r\sin b\sin(B+t)\\
z&=r\sin c\sin(C+t)
\end{aligned}
\]

Nun bleibt noch, die Größen $a$, $A$ usw., die allein von der Lage der
Planetenbahn und nicht von der jeweiligen Position des Planeten in ihr
abhängen, aus der Neigung der Bahnebene und der Länge des aufsteigenden Knotens
abzuleiten. Wir bezeichnen die Neigung mit $i$ und die Länge des Knotens mit
$\AscendingNode$. Aus dem Dreieck $\mathfrak{x}pP$ ergeben sich die folgenden
drei Gleichungen
\[
\begin{aligned}
\operatorname{cotang}\mathfrak{x}pP&=\frac{\sin pP\operatorname{cotang}\mathfrak{x}P-\cos pP\cos pP\mathfrak{x}}{\sin pP\mathfrak{x}}\\
\cos\mathfrak{x}p&=\cos pP\cos\mathfrak{x}P+\sin pP\sin\mathfrak{x}P\cos pP\mathfrak{x}\\
\sin\mathfrak{x}p&=\frac{\sin\mathfrak{x}P\sin pP\mathfrak{x}}{\sin\mathfrak{x}pP}
\end{aligned}
\]

Ebenso liefern die Dreiecke $\mathfrak{y}pP$, $\mathfrak{z}pP$ jeweils drei
ähnliche Gleichungen. Sie brauchen hier nicht ausgeschrieben zu werden, da
man in den drei obigen lediglich $\mathfrak{x}$ durch $\mathfrak{y}$ bzw.
$\mathfrak{z}$ ersetzen muss. Nun gilt
\[
pP=i,\quad pP\mathfrak{x}=90^0-\AscendingNode,\quad pP\mathfrak{y}=pP\mathfrak{z}=180^0-\AscendingNode
\]
Mit diesen und den übrigen Einsetzungen erhalten wir für unsere neun Gleichungen
\[
\begin{aligned}
\operatorname{cotang}A&=-\cos i\operatorname{tang}\AscendingNode\\
\cos a&=\sin i\sin\AscendingNode\\
\sin a&=\frac{\cos\AscendingNode}{\sin A}
\end{aligned}
\]
\SourceLeaf{104}%
\[
\begin{aligned}
\operatorname{cotang}B&=\frac{-\sin i\operatorname{tang}\varepsilon+\cos i\cos\AscendingNode}{\sin\AscendingNode}\\
\cos b&=-\cos i\sin\varepsilon-\sin i\cos\varepsilon\cos\AscendingNode\\
\sin b&=\frac{\cos\varepsilon\sin\AscendingNode}{\sin B}\\
\operatorname{cotang}C&=\frac{\sin i\operatorname{cotang}\varepsilon+\cos i\cos\AscendingNode}{\sin\AscendingNode}\\
\cos c&=\cos i\cos\varepsilon-\sin i\sin\varepsilon\cos\AscendingNode\\
\sin c&=\frac{\sin\varepsilon\sin\AscendingNode}{\sin C}
\end{aligned}
\]

Ob $A$, $B$ und $C$ in den ersten beiden oder den letzten beiden
Quadranten anzusetzen sind, entscheidet man so, dass die Sinuswerte von
$a$, $b$ und $c$ positiv sind. Man setzt also $A$ in den ersten beiden
Quadranten an, wenn $\cos\AscendingNode$ positiv ist, und $B$ und $C$ in
denselben, wenn $\sin\AscendingNode$ positiv ist. In den entgegengesetzten
Fällen setzt man sie in den letzten beiden Quadranten an. $B$ und $C$
liegen also in den ersten oder letzten beiden Quadranten, je nachdem,
ob $\AscendingNode$ in den ersten oder den letzten beiden liegt.

Die vierte, fünfte, siebte und achte dieser Gleichungen lassen sich durch
die Einführung von Hilfswinkeln noch bequemer gestalten. Das ist auf zwei
Arten möglich.

\textit{Erstens}. Setzt man $\frac{\operatorname{tang}i}{\cos\AscendingNode}=\operatorname{tang}E$
und $\operatorname{tang}i\cos\AscendingNode=\operatorname{tang}F$, so ergibt sich
\[
\begin{aligned}
\operatorname{cotang}B&=\frac{\sin i\cos(E+\varepsilon)}{\sin\AscendingNode\cos\varepsilon\sin E}=\frac{\cos i\cos(E+\varepsilon)}{\operatorname{tang}\AscendingNode\cos\varepsilon\cos E}\\
\cos b&=-\frac{\cos i\sin(F+\varepsilon)}{\cos F}=-\frac{\sin i\cos\AscendingNode\sin(F+\varepsilon)}{\sin F}\\
\operatorname{cotang}C&=\frac{\sin i\sin(E+\varepsilon)}{\sin\AscendingNode\sin\varepsilon\sin E}=\frac{\cos i\sin(E+\varepsilon)}{\operatorname{tang}\AscendingNode\sin\varepsilon\cos E}\\
\cos c&=\frac{\cos i\cos(F+\varepsilon)}{\cos F}=\frac{\sin i\cos\AscendingNode\cos(F+\varepsilon)}{\sin F}
\end{aligned}
\]

\textit{Zweitens}. Setzt man $\frac{\operatorname{tang}\varepsilon}{\cos\AscendingNode}=\operatorname{tang}G$
und $\operatorname{tang}\varepsilon\cos\AscendingNode=\operatorname{tang}H$, so ergibt sich
\[
\begin{aligned}
\operatorname{cotang}B&=\frac{\cos(G+i)}{\operatorname{tang}\AscendingNode\cos G}=\frac{\operatorname{tang}\varepsilon\cos(G+i)}{\sin\AscendingNode\sin G}\\
\cos b&=-\frac{\sin\varepsilon\sin(G+i)}{\sin G}=-\frac{\cos\AscendingNode\cos\varepsilon\sin(G+i)}{\cos G}\\
\operatorname{cotang}C&=\frac{\sin(H+i)}{\sin\AscendingNode\operatorname{tang}\varepsilon\cos H}=\frac{\sin(H+i)}{\operatorname{tang}\AscendingNode\sin H}\\
\cos c&=\frac{\cos\varepsilon\cos(H+i)}{\cos H}=\frac{\sin\varepsilon\cos\AscendingNode\cos(H+i)}{\sin H}
\end{aligned}
\]

\SourceLeaf{105}%
Es dürfte sich lohnen, noch einige Beziehungen zwischen den Größen
$A$, $a$ usw. herzuleiten. Aus dem sphärischen Dreieck $\mathfrak{x}p\mathfrak{y}$
ergibt sich $\cos\mathfrak{x}\mathfrak{y}=\cos\mathfrak{x}p\cos\mathfrak{y}p+\sin\mathfrak{x}p\sin\mathfrak{y}p\cos\mathfrak{x}p\mathfrak{y}$.
Nun ist $\mathfrak{x}\mathfrak{y}=90^0$ und
$\mathfrak{x}p\mathfrak{y}=\mathfrak{x}pP-\mathfrak{y}pP=A-B$. Also gilt
\[
\cos(A-B)=-\operatorname{cotang}a\operatorname{cotang}b
\]
Ebenso ergeben die Dreiecke $\mathfrak{y}p\mathfrak{z}$, $\mathfrak{z}p\mathfrak{x}$
\[
\begin{aligned}
\cos(B-C)&=-\operatorname{cotang}b\operatorname{cotang}c\\
\cos(C-A)&=-\operatorname{cotang}c\operatorname{cotang}a
\end{aligned}
\]

Ferner gilt im Dreieck $\mathfrak{x}p\mathfrak{y}$ die Gleichung
$\cos a=\cos p\mathfrak{y}\mathfrak{x}\sin b$ und im Dreieck
$\mathfrak{y}p\mathfrak{z}$ die Gleichung $\sin\mathfrak{z}\mathfrak{y}p=\sin\mathfrak{y}p\mathfrak{z}\sin c$.
Da nun $\mathfrak{z}\mathfrak{y}p+p\mathfrak{y}\mathfrak{x}=\mathfrak{z}\mathfrak{y}\mathfrak{x}=90^0$ ist,
gilt $\cos a=\sin b\sin c\sin\mathfrak{y}p\mathfrak{z}$ oder, da
$\mathfrak{y}p\mathfrak{z}=B-C$ ist,
\[
\sin(B-C)=\frac{\cos a}{\sin b\sin c}
\]
Ganz ähnlich erhält man
\[
\begin{aligned}
\sin(C-A)&=\frac{\cos b}{\sin c\sin a}\\
\sin(A-B)&=\frac{\cos c}{\sin a\sin b}
\end{aligned}
\]

Verbindet man diese Gleichungen mit den vorherigen, so ergibt sich noch
\[
\begin{aligned}
\operatorname{cotang}(A-B)&=-\frac{\cos a\cos b}{\cos c}=\operatorname{tang}(F+\varepsilon)\cos a\\
\operatorname{cotang}(B-C)&=-\frac{\cos b\cos c}{\cos a}\\
\operatorname{cotang}(C-A)&=-\frac{\cos c\cos a}{\cos b}=\operatorname{cotang}(F+\varepsilon)\cos a
\end{aligned}
\]
\[
\begin{aligned}
\cos a^2&=\operatorname{cotang}(A-B)\operatorname{cotang}(C-A)\\
\cos b^2&=\operatorname{cotang}(B-C)\operatorname{cotang}(A-B)\\
\cos c^2&=\operatorname{cotang}(C-A)\operatorname{cotang}(B-C)
\end{aligned}
\]
und auf ähnliche Weise lassen sich die Quadrate der Sinus- und Tangenswerte
der Seiten $a$, $b$, $c$ durch die Winkel $A-B$, $B-C$, $C-A$ ausdrücken.

\SourceLeaf{106}%
Um den Gebrauch dieser Formeln zu erläutern, wenden wir einige davon auf
Pallas an und legen dabei die neuesten [VII.] Elemente dieses Planeten für
1803 zugrunde. Wir setzen also
\[
\begin{aligned}
i&=34^0\;38'\;1''1\\
\AscendingNode&=172\;28\;13.7\\
\varepsilon&=23\;27\;55.8\quad\text{(mittlere Schiefe nach \textsc{Maskelyne} für 1803).}
\end{aligned}
\]

Mit diesen Elementen ergibt sich die folgende Rechnung. Das n neben
einem Logarithmus zeigt an, dass die zugehörige Größe negativ ist.

\begin{center}\begin{minipage}{\textwidth}\centering
\begin{tabular}{@{}lcl@{}}
$\log\cos i$ & $ .\quad .\quad .$ & $9.9152958$\\
$\log\operatorname{tang}\AscendingNode$ & $ .\quad .$ & $9.1211553\,\mathrm{n}$\\\hline
$\log\operatorname{cotang}A$ & $.$ & $9.0364511$
\end{tabular}\par\smallskip
 also $A=263^0\;47'\;35''4$
\end{minipage}\end{center}
\begin{center}\begin{minipage}{\textwidth}\centering
\begin{tabular}{@{}lcl@{}}
$\log\cos\AscendingNode$ & $ .\quad .\quad .$ & $9.9962390\,\mathrm{n}$\\
$\log\sin A$ & $ .\quad .\quad .$ & $9.9974467\,\mathrm{n}$\\\hline
$\log\sin a$ & $ .\quad .\quad .$ & $9.9987923$
\end{tabular}
\end{minipage}\end{center}
\begin{center}\begin{minipage}{\textwidth}\centering
\begin{tabular}{@{}lcl@{}}
$\log\sin i$ & $ .\quad .\quad .$ & $9.7545982$\\
$\log\sin\AscendingNode$ & $ .\quad .\quad .$ & $9.1173944$\\\hline
$\log\cos a$ & $ .\quad .\quad .$ & $8.8719926$
\end{tabular}\par\smallskip
 hieraus $a=85^0\;43'\;44''8$
\end{minipage}\end{center}
\begin{center}\begin{minipage}{\textwidth}\centering
\begin{tabular}{@{}lcl@{}}
$\log\operatorname{tang}i$ & $ .\quad .$ & $9.8393024$\\
$\log\cos\AscendingNode$ & $ .\quad .\quad .$ & $9.9962390\,\mathrm{n}$\quad also\\\hline
\end{tabular}
\[
\begin{gathered}
\log\operatorname{tang}E\quad .\quad .\quad 9.8430634\,\mathrm{n},\\
E=145^0\;8'\;2''4;\quad E+\varepsilon=168^0\;35'\;58''2\\
\log\operatorname{tang}F\quad .\quad .\quad 9.8355414\,\mathrm{n},\\
F=145\;35\;52.9;\quad F+\varepsilon=169\;3\;48.7
\end{gathered}
\]
\end{minipage}\end{center}
\begin{center}\begin{minipage}{\textwidth}\centering
\begin{tabular}{@{}lcl@{}}
$\log\cos i$ & $ .\quad .\quad .$ & $9.9152958$\\
Compl. $\log\operatorname{tang}\AscendingNode$ & & $0.8788447\,\mathrm{n}$\\
Compl. $\log\cos E$ & & $0.0859260\,\mathrm{n}$\\\hline
$\log\mathrm{const}.$ & $ .\quad .$ & $0.8800665$
\end{tabular}
\end{minipage}\end{center}
\begin{center}\begin{minipage}{\textwidth}\centering
\begin{tabular}{@{}lcl@{}}
$\log\cos(E+\varepsilon)$ & $.$ & $9.9913455\,\mathrm{n}$\\
Compl. $\log\cos\varepsilon$ & & $0.0374886$\\\hline
$\log\operatorname{cotang}B$ & $.$ & $0.9089006\,\mathrm{n}$
\end{tabular}\par\smallskip
 hieraus $B=172^0\;58'\;7''4$,
\end{minipage}\end{center}
\SourceLeaf{107}%
\begin{center}\begin{minipage}{\textwidth}\centering
\begin{tabular}{@{}lcl@{}}
$\log\mathrm{const}.$ & $ .\quad .$ & $0.8800665$\\
$\log\sin(E+\varepsilon)$ & $.$ & $9.2959318$\\
Compl. $\log\sin\varepsilon$ & & $0.3999023$\\\hline
$\log\operatorname{cotang}C$ & $.$ & $0.5759006$
\end{tabular}\par\smallskip
 $C=14^0\;52'\;12''5$
\end{minipage}\end{center}
\begin{center}\begin{minipage}{\textwidth}\centering
\begin{tabular}{@{}lcl@{}}
$\log\cos\varepsilon$ & $ .\quad .\quad .$ & $9.9625114$\\
$\log\sin\AscendingNode$ & $ .\quad .\quad .$ & $9.1173944$\\
Compl. $\log\sin B$ & & $0.9121791$\\\hline
$\log\sin b$ & $ .\quad .\quad .$ & $9.9920849$
\end{tabular}
\end{minipage}\end{center}
\begin{center}\begin{minipage}{\textwidth}\centering
\begin{tabular}{@{}lcl@{}}
$\log\sin\varepsilon$ & $ .\quad .\quad .$ & $9.6000977$\\
$\log\sin n$ & $ .\quad .\quad .$ & $9.1173944$\\
Compl. $\log\sin C$ & & $0.5906942$\\\hline
$\log\sin c$ & $ .\quad .\quad .$ & $9.3081863$
\end{tabular}
\end{minipage}\end{center}
\begin{center}\begin{minipage}{\textwidth}\centering
\begin{tabular}{@{}lcl@{}}
$\log\cos i$ & $ .\quad .\quad .$ & $9.9152958$\\
$\log\cos F$ & $ .\quad .$ & $9.9165035\,\mathrm{n}$\\\hline
 & & $9.9987923\,\mathrm{n}$
\end{tabular}
\end{minipage}\end{center}
\begin{center}\begin{minipage}{\textwidth}\centering
\begin{tabular}{@{}lcl@{}}
$\log\sin(F+\varepsilon)$ & $.$ & $9.2781142$\\
$\log\cos(F+\varepsilon)$ & $.$ & $9.9920399\,\mathrm{n}$\\\hline
$\log\cos b$ & $ .\quad .\quad .$ & $9.2769065$\\
$\log\cos c$ & $ .\quad .\quad .$ & $9.9908322$
\end{tabular}\par\smallskip
 also
$\begin{aligned}b&=79^0\;5'\;39''4\\c&=11\;43\;52.8\end{aligned}$
\end{minipage}\end{center}

Braucht man nur die Sinuswerte von $a$, $b$, $c$, so ist die Berechnung
ihrer Kosinuswerte entbehrlich und damit auch der Hilfswinkel $F$. Will man
aber auch $a$, $b$, $c$ selbst kennen, so helfen die Kosinuswerte, die
später noch einmal gebraucht werden, die durch die Sinuswerte allein
verbleibenden Mehrdeutigkeiten aufzulösen. Liegen die Sinuswerte näher
bei 1, so ermöglichen die Kosinuswerte eine genauere Bestimmung und zugleich
eine Kontrolle der Rechnung. Dafür ist auch die Gleichung
$\frac{\cos i}{\cos F}=\pm\sin a$ nützlich. Das obere Vorzeichen gilt, wenn
$F$ und $A$ beide in den ersten oder beide in den letzten zwei Quadranten
liegen. Das untere gilt, wenn $F$ in der anderen Hälfte des Kreisumfangs
als $A$ angesetzt wird.
\SourceLeaf{108}%
(Der Grund dafür ergibt sich aus der Bemerkung, dass $F$ im ersten Fall
mit dem Winkel $P\mathfrak{x}p$ übereinstimmt und im zweiten um $180^0$ davon
abweicht.)

Die Größen $\varepsilon$, $\AscendingNode$, $i$ unterliegen säkularen
Änderungen. Dasselbe gilt daher für die davon abhängigen Größen $A$, $a$,
$B$, $b$, $C$, $c$. Sind die jährlichen Änderungen der erstgenannten Größen
bekannt, so lassen sich die Änderungen von $A$, $a$ usw. mit leicht
herzuleitenden Differentialformeln berechnen, bei denen wir uns hier nicht
aufhalten wollen. Man kann auch die Werte von $A$, $a$ usw. für eine weiter
entfernte Epoche neu berechnen und daraus ihre jährlichen Änderungen ableiten.

Außerdem erfahren diese Größen durch die Nutation periodische Änderungen,
die mit jedem Umlauf der Mondknoten wiederkehren. Da die geozentrische Position
des Planeten in Bezug auf den wahren Äquator gesucht ist, muss man eigentlich
für $\varepsilon$ die wahre Schiefe der Ekliptik statt der mittleren und für
$\AscendingNode$ den Abstand des aufsteigenden Knotens vom wahren statt vom
mittleren Äquinoktialpunkt nehmen. Die daraus entstehenden periodischen
Änderungen lassen sich mit denselben Differentialformeln wie die säkularen
Änderungen berechnen und in eine Tabelle eintragen, deren Argument die Länge
des Mondknotens ist. Sind zahlreiche geozentrische Positionen für einen nicht
zu langen Zeitraum zu berechnen und fehlt eine solche Tabelle, so berechnet
man am bequemsten die wahren Werte von $A$, $a$ usw. für zwei Epochen zu Anfang
und Ende des Zeitraums unmittelbar aus den wahren Werten von $\varepsilon$,
$i$, $\AscendingNode$ und gewinnt daraus die Werte für die dazwischenliegenden
Zeiten durch einfache Interpolation. Über ein Jahr kann man diese Änderungen
ohne Bedenken als gleichförmig ansehen.

Man könnte auch die durch die Nutation bedingten periodischen Änderungen
ganz weglassen und die mittleren Werte von $A$, $a$ usw. verwenden. Dann
müsste man aber auch bei der Erde für $\varepsilon$ die mittlere Schiefe der
Ekliptik verwenden und bei der Länge $\lambda$ die Nutation weglassen, um
den Abstand vom mittleren Äquinoktium zu erhalten. So erhält man die
geozentrische Rektaszension und Deklination des Planeten in Bezug auf den
\textit{mittleren} Äquator. Daraus leitet man seine Position zum \textit{wahren}
Äquator ebenso ab, wie man durch Berücksichtigung der Nutation die mittlere
Position eines Fixsterns auf seine scheinbare Position reduziert.

\SourceLeaf{109}%
Wir haben jetzt nur noch einige Bemerkungen über die Störungen hinzuzufügen.
Die Breitenstörungen, um die es hier allein geht, sind bei allen älteren
Planeten so unbedeutend, dass man sie mit Recht ganz vernachlässigen kann.
Nur bei Ceres und Pallas müssen sie wegen der starken Neigung ihrer Bahnen
zur Jupiterbahn berücksichtigt werden. Das kann auf zwei Wegen geschehen.
Man kann die Elemente, die die Lage der Bahn bestimmen, also die Neigung
und die Länge des Knotens, als veränderlich betrachten und ihre mittleren
Werte durch periodische Gleichungen berichtigen. Oder man untersucht
unmittelbar, wie weit der Planet durch fremde Kräfte aus der mittleren
Ebene seiner Bahn herausgedrängt wird. Im ersten Fall überträgt man die
Änderungen auch auf die Größen $A$, $a$ usw. und fügt zu den von der Nutation
abhängigen periodischen Gleichungen weitere hinzu. Deren Argumente stimmen
mit denen der Gleichungen für die Neigung und die Länge des Knotens überein.
Dieses Verfahren war jedoch bisher nicht üblich. Bei der zweiten Methode
dagegen geben die Störungstafeln die Störung der heliozentrischen Breite an,
also die heliozentrische Breite des Planeten über der mittleren Ebene seiner
Bahn. Sie sei $=\ObservedBeta$ und werde zum Nordpol hin als positiv, zum
Südpol hin als negativ gerechnet. Im sphärischen Dreieck $\mathfrak{x}pk$
ist die Seite $pk$ also $=90^0$ im vorigen Fall und $=90^0-\ObservedBeta$ hier.
Daraus folgt
\[
x=r\cos\mathfrak{x}k=r(\sin\ObservedBeta\cos a+\cos\ObservedBeta\sin a\sin(t+A))
\]
und ebenso
\[
\begin{aligned}
y&=r(\sin\ObservedBeta\cos b+\cos\ObservedBeta\sin b\sin(t+B))\\
z&=r(\sin\ObservedBeta\cos c+\cos\ObservedBeta\sin c\sin(t+C))
\end{aligned}
\]

Da $\ObservedBeta$ hier höchstens einige Minuten betragen kann, darf man
$\cos\ObservedBeta=1$ und $\sin\ObservedBeta=\ObservedBeta$ setzen. Wegen der
Störungen braucht man also zu den ohne sie gefundenen Werten von $x$, $y$,
$z$ nur die Größen $\ObservedBeta r\cos a$, $\ObservedBeta r\cos b$,
$\ObservedBeta r\cos c$ zu addieren. Dabei muss $\ObservedBeta$ in Teilen
des Radius ausgedrückt werden.

\medskip
[Handschriftliche Bemerkungen zu S.\,100.]
\[
\sin a^2=-\frac{\cos(B-C)}{\sin(A-B)\sin(C-A)}
\]
\SourceLeaf{110}%
\[
\operatorname{tang}a^2=-\frac{\cos(B-C)}{\cos(A-B)\cos(C-A)}
\]
\[
\cos a^2+\cos b^2+\cos c^2=1
\]
\[
\begin{aligned}
-\frac{\cos a}{\cos b}&=\operatorname{tang}(B-C)\cos=\frac{\operatorname{cotang}(C-A)}{\cos c}\\
&=\operatorname{tang}\AscendingNode'\ \text{in plano ipsius }Z
\end{aligned}
\]
\[
\begin{aligned}
\sin\AscendingNode'&=\sin b\sin(B-C)=\frac{\cos a}{\sin c}\\
\cos\AscendingNode'&=-\frac{\cos b}{\sin c}
\end{aligned}
\]

\begin{center}Differentialänderungen:\end{center}
\[
2\mathrm{d}a=-\frac{\mathrm{d}A-\mathrm{d}B}{\sin(A-B)\operatorname{cotang}b}-\frac{\mathrm{d}C-\mathrm{d}A}{\sin(C-A)\operatorname{cotang}c}
\]

Umgekehrt lassen sich aus $a$, $b$, $c$ die Größen $\AscendingNode$
und $i$ mit den folgenden Formeln bestimmen
\[
\begin{aligned}
\cos a&=\sin i\sin\AscendingNode\\
-\cos b\cos\varepsilon-\cos c\sin\varepsilon&=\sin i\cos\AscendingNode\\
-\cos b\sin\varepsilon+\cos c\cos\varepsilon&=\cos i
\end{aligned}
\]
\[
\begin{aligned}
\sin a&=k\\
\cos b&=k\cos\theta\\
\cos c&=k\sin\theta
\end{aligned}
\]
\[
\begin{aligned}
\cos a&=\sin i\sin\AscendingNode\\
-k\cos(\theta-\varepsilon)&=\sin i\cos\AscendingNode\\
k\sin(\theta-\varepsilon)&=\cos i
\end{aligned}
\]
\begin{center}\rule{.25\linewidth}{.4pt}\end{center}
